\documentclass[10pt,a4paper]{amsart}
\usepackage[margin=19mm]{geometry}
\usepackage{amsmath,amssymb,mathtools}
\usepackage[scr=rsfs,cal=euler]{mathalfa}
\usepackage{xfrac}
\usepackage[unicode,hidelinks,bookmarks=false]{hyperref}
\newcommand{\ie}{i.e.\ }
\newcommand{\cf}{cf.\ }
\newcommand{\resp}{respectively\ }
\newcommand{\Subsection}{\subsection}
\providecommand{\nobreakdash}{\nobreak\hbox{-}}
\setlength{\emergencystretch}{2em}
\multlinegap0pt
\usepackage{amssymb}
\usepackage{enumerate}
\usepackage[shortlabels]{enumitem}
\usepackage{xcolor}
\usepackage{mathtools}
\newtheorem{proposition}{Proposition}
\newcommand{\RR}{\mathbb{R}}
\title{Quantitative limit theorems for generalized P\'olya urns with applications to random tree models}
\author{Morenikeji Neri${}^{\MakeLowercase a}$ and Pedro Pinto${}^{\MakeLowercase b}$}
\date{Source: arXiv:2607.07309v1 (CC BY 4.0)}
\begin{document}
\maketitle

\noindent\emph{Source and licence.} Quantitative limit theorems for generalized P\'olya urns with applications to random tree models. Version arXiv:2607.07309v1 (CC BY 4.0), distributed by arXiv under the Creative Commons Attribution 4.0 International licence, http://creativecommons.org/licenses/by/4.0/, as recorded in the arXiv licence metadata for this exact version and shown on the abstract page https://arxiv.org/abs/2607.07309v1. Full licence notice: \texttt{licenses/arxiv-2607.07309-CC-BY-4.0.txt}.\par\smallskip

\noindent\textit{Editorial scope.} Counted unit: the article's proposition prop:urn:SA (source0.tex lines 549--562). The article's complete original proof of it is delivered with it (source0.tex lines 564--596, one \texttt{proof} environment of 33 lines). No mathematics is written, repaired or re-derived: the statement and the proof are the article's own lines, in the article's own notation, and no retained line is substituted or re-flowed.  No further source-local context is needed: every cross-reference of every retained line resolves inside the two delivered spans.  Nothing was omitted from any retained span, so the package carries no omission record.  No retained line contains a citation, so the wrapper carries no bibliography and no bibliography entry.  No retained line contains a figure environment, an image-inclusion command or drawing code, so no image is delivered and the package declares no asset dependency.  Editorial formatting only, in addition: an AMS \texttt{amsart} preamble that restates the article's own declarations and macros used by the retained lines, namely the environments \texttt{proposition} on separate counters, together with the standard packages those lines need.  The retained environments print in this document as Proposition 1 in order of appearance, whereas the article prints the same items with the numbering of its own full document.  The complete native source of the article as distributed in the arXiv e-print for this exact version is bundled unchanged as \texttt{sources/204695/source0.tex}.
\begin{proposition}\label{prop:urn:SA}
For all $n \geq 0$,
\[
    \hat{U}(n+1) = \hat{U}(n) + \gamma_n\bigl(F(\hat{U}(n)) + \Delta M_{n+1}\bigr),
\]
where $\gamma_n = \frac{1}{T(n+1)}$, $ \Delta M_{n+1}
    := R_{\xi(n+1)} 
      - \mathbb{E}\left[R_{\xi(n+1)} \mid \mathcal{F}_n\right]$, and  $F: \RR^d \to \RR^d$ is a function given by $F(x) = \sum_{i=0}^{d-1} x_i\bigl(R_i - \lambda\, x\bigr)$. Moreover, on the simplex
\[
    \mathcal{S} = \Bigl\{(x_0, \ldots, x_{d-1}) \in [0,1]^d : \sum_{i=0}^{d-1} x_i = 1\Bigr\},
\]
we have $F(x)= x(R-\lambda I)$ which implies $F(\hat{U}(n)) = \hat{U}(n)(R-\lambda I)$.

\end{proposition}

\begin{proof}
From the recursive definition of $(U(n))_{n \geq 0}$, we have
\begin{align*}
    \hat{U}(n+1)
    &= \frac{U(n) +R_{\xi(n+1)}}{T(n+1)}
     = \frac{U(n)}{T(n)} \cdot \frac{T(n)}{T(n+1)} + \frac{R_{\xi(n+1)}}{T(n+1)} \\[6pt]
    &= \hat{U}(n) \cdot \left(1 - \frac{\lambda}{T(n+1)}\right)
       + \frac{R_{\xi(n+1)}}{T(n+1)}
       =\hat{U}(n)+ \frac{R_{\xi(n+1)} - \lambda\,\hat{U}(n)}{T(n+1)}.
\end{align*}
Now, we set
\[
    \Delta M_{n+1}
    = R_{\xi(n+1)} - \lambda\,\hat{U}(n)
      - \mathbb{E}\!\left[R_{\xi(n+1)} - \lambda\,\hat{U}(n)\mid \mathcal{F}_n\right]=R_{\xi(n+1)} 
      - \mathbb{E}\!\left[R_{\xi(n+1)} \mid \mathcal{F}_n\right].
\]
Since $\mathbb{P}(\xi(n+1) = i \mid \mathcal{F}_n) = \hat{U}_i(n)$, we get that
\[
    \mathbb{E}\!\left[R_{\xi(n+1)} - \lambda\,\hat{U}(n)
    \mid \mathcal{F}_n\right]
    = \sum_{i=1}^{d} \hat{U}_i(n)\bigl(R_i - \lambda\,\hat{U}(n)\bigr)=F(\hat{U}(n)).
\]
It thus follows that,
\[
    \hat{U}(n+1) = \hat{U}(n) + \frac{1}{T(n+1)}\bigl(F(\hat{U}(n)) + \Delta M_{n+1}\bigr).
\]
For $x \in \mathcal{S}$, we easily see for all $0\leq j\leq d-1$
\begin{align*}
    F(x)_j&= \sum_{i=0}^{d-1}x_iR_{ij}-\lambda x_j\sum_{i=0}^{d-1}x_i\\
    &=\sum_{i=0}^{d-1}x_iR_{ij}-\lambda x_j=(x(R-\lambda I)_j.
\end{align*}
\end{proof}

\end{document}
